\documentclass[11pt,fontset=fandol]{ctexart}
\usepackage[a4paper,margin=23mm]{geometry}
\usepackage{amsmath,amssymb,mathtools,bm,mathrsfs,esint}
\usepackage{graphicx,longtable,booktabs,array,calc}
\usepackage[normalem]{ulem}
\usepackage{xcolor,hyperref}
\hypersetup{hidelinks}
\setcounter{secnumdepth}{-1}
\newcounter{none}
\providecommand{\tightlist}{\setlength{\itemsep}{0pt}\setlength{\parskip}{0pt}}
\providecommand{\passthrough}[1]{#1}
\usepackage{amssymb,mathrsfs,esint,bm,newunicodechar}
\allowdisplaybreaks
\setlength{\emergencystretch}{3em}
\providecommand{\llangle}{\langle\!\langle}
\providecommand{\rrangle}{\rangle\!\rangle}
\newunicodechar{∞}{\ensuremath{\infty}}
\newunicodechar{∈}{\ensuremath{\in}}
\newunicodechar{∑}{\ensuremath{\sum}}
\newunicodechar{∫}{\ensuremath{\int}}
\newunicodechar{≤}{\ensuremath{\leq}}
\newunicodechar{≥}{\ensuremath{\geq}}
\newunicodechar{≠}{\ensuremath{\neq}}
\newunicodechar{→}{\ensuremath{\to}}
\newunicodechar{⇒}{\ensuremath{\Rightarrow}}
\newunicodechar{α}{\ensuremath{\alpha}}
\newunicodechar{β}{\ensuremath{\beta}}
\newunicodechar{γ}{\ensuremath{\gamma}}
\newunicodechar{λ}{\ensuremath{\lambda}}
\newunicodechar{φ}{\ensuremath{\varphi}}
\newunicodechar{ε}{\ensuremath{\varepsilon}}
\newunicodechar{σ}{\ensuremath{\sigma}}
\newunicodechar{θ}{\ensuremath{\theta}}
\newunicodechar{δ}{\ensuremath{\delta}}
\newunicodechar{η}{\ensuremath{\eta}}
\newunicodechar{π}{\ensuremath{\pi}}
\newunicodechar{ξ}{\ensuremath{\xi}}
\newunicodechar{Σ}{\ensuremath{\Sigma}}
\newunicodechar{Δ}{\ensuremath{\Delta}}
\newunicodechar{∂}{\ensuremath{\partial}}
\newunicodechar{∀}{\ensuremath{\forall}}
\newunicodechar{∃}{\ensuremath{\exists}}
\newunicodechar{√}{\ensuremath{\surd}}
\newunicodechar{−}{\ensuremath{-}}
\newunicodechar{μ}{\ensuremath{\mu}}
\newunicodechar{ρ}{\ensuremath{\rho}}
\newunicodechar{τ}{\ensuremath{\tau}}
\newunicodechar{ω}{\ensuremath{\omega}}
\newunicodechar{Ω}{\ensuremath{\Omega}}
\newunicodechar{χ}{\ensuremath{\chi}}
\newunicodechar{ψ}{\ensuremath{\psi}}
\newunicodechar{Γ}{\ensuremath{\Gamma}}
\newunicodechar{Φ}{\ensuremath{\Phi}}
\newunicodechar{Λ}{\ensuremath{\Lambda}}
\newunicodechar{ℝ}{\ensuremath{\mathbb{R}}}
\newunicodechar{ℤ}{\ensuremath{\mathbb{Z}}}
\newunicodechar{ℕ}{\ensuremath{\mathbb{N}}}
\newunicodechar{△}{\ensuremath{\triangle}}
\newunicodechar{⊗}{\ensuremath{\otimes}}
\begin{document}
\section{2022年高等代数参考答案}\label{ux5e74ux9ad8ux7b49ux4ee3ux6570ux53c2ux8003ux7b54ux6848}

来源：\texttt{HDU/601数学分析/2022/2022年杭州电子科技大学数分高代真题及答案(双科合卷).pdf}。原
PDF 页序：16, 17, 18, 19, 20, 21, 22, 23, 24。

\section{杭州电子科技大学2022年高等代数真题解析}\label{ux676dux5ddeux7535ux5b50ux79d1ux6280ux5927ux5b662022ux5e74ux9ad8ux7b49ux4ee3ux6570ux771fux9898ux89e3ux6790}

\subsection{一、计算行列式}\label{ux4e00ux8ba1ux7b97ux884cux5217ux5f0f}

\textbf{依据题面重建的补充推导：}原解析此处的中间矩阵无法可靠识读，以下推导独立恢复对应结果。

将原矩阵的列逆序，记所得行列式为 \(H_n\)，有 \[
D_n=(-1)^{n(n-1)/2}H_n,\qquad
H_n=\begin{vmatrix}
a_1&x&\cdots&x\\
x&a_2&\cdots&x\\
\vdots&\vdots&\ddots&\vdots\\
x&x&\cdots&a_n
\end{vmatrix}.
\] 原解析以 \(x\ne a_i\)（\(i=2,\ldots,n\)）开始消元；可先在更方便的
\(x\ne a_i\)（全部 \(i\)）下进行同一计算，写成对角矩阵加秩一矩阵： \[
H_n=\det\!\left(\operatorname{diag}(a_1-x,\ldots,a_n-x)
+x\boldsymbol1\boldsymbol1^T\right),\qquad
\boldsymbol1=(1,\ldots,1)^T.
\] 由秩一行列式公式， \[
H_n=\prod_{i=1}^n(a_i-x)\left(1+x\sum_{i=1}^n\frac1{a_i-x}\right)
=\prod_{i=1}^n(a_i-x)+x\sum_{i=1}^n\prod_{j\ne i}(a_j-x).
\] 最后一式是多项式恒等式，因此也覆盖某些 \(x=a_i\) 的情况。答案为 \[
\boxed{D_n=(-1)^{n(n-1)/2}\left[\prod_{i=1}^n(a_i-x)+x\sum_{i=1}^n\prod_{j\ne i}(a_j-x)\right]}.
\]

\subsection{二、多项式互素}\label{ux4e8cux591aux9879ux5f0fux4e92ux7d20}

由 \(ad-bc\ne0\)，\(a,b\) 不全为零。

当 \(a=0\) 时，\(b,c\ne0\)，于是 \[
\begin{aligned}
(af+bg,cf+dg)&=(bg,cf+dg)\\
&=\left(bg,cf+dg-\frac dbbg\right)\\
&=(g,cf)=(f,g).
\end{aligned}
\] 当 \(a\ne0\) 时， \[
\begin{aligned}
(af+bg,cf+dg)
&=\left(af+bg,cf+dg-\frac ca(af+bg)\right)\\
&=\left(af+bg,\frac{ad-bc}{a}g\right)\\
&=(af+bg,g)=(af,g)=(f,g).
\end{aligned}
\] 所以两个首一最大公因式相同，题中充要条件成立。

\subsection{三、齐次线性方程组}\label{ux4e09ux9f50ux6b21ux7ebfux6027ux65b9ux7a0bux7ec4}

令 \(s=a_1+a_2+a_3+a_4\ne0\)。系数矩阵为 \[
A=\begin{pmatrix}
a_1+b&a_2&a_3&a_4\\a_1&a_2+b&a_3&a_4\\
a_1&a_2&a_3+b&a_4\\a_1&a_2&a_3&a_4+b
\end{pmatrix}
=bE+\begin{pmatrix}1\\1\\1\\1\end{pmatrix}(a_1,a_2,a_3,a_4).
\] 由秩一矩阵的性质，特征值为 \(b,b,b,b+s\)。这里 \[
(a_1,a_2,a_3,a_4)\begin{pmatrix}1\\1\\1\\1\end{pmatrix}=s.
\] 特征值 \(b\) 的特征向量是 \((a_1,a_2,a_3,a_4)x=0\) 的非零解。至少一个
\(a_i\ne0\)；不妨设 \(a_4\ne0\)，则通解为 \[
x=\begin{pmatrix}x_1\\x_2\\x_3\\-(a_1x_1+a_2x_2+a_3x_3)/a_4\end{pmatrix}
=x_1\begin{pmatrix}1\\0\\0\\-a_1/a_4\end{pmatrix}
+x_2\begin{pmatrix}0\\1\\0\\-a_2/a_4\end{pmatrix}
+x_3\begin{pmatrix}0\\0\\1\\-a_3/a_4\end{pmatrix}.
\] 特征值 \(b+s\) 的特征向量可取
\((1,1,1,1)^T\)。原OCR将该列向量误识别为五行，此处按四维题面恢复为四行。

（1）当且仅当 \(b\ne0\) 且 \(b\ne-s\) 时，\(A\) 可逆，方程组只有零解。

（2）\(b=0\) 时，上式中的三个列向量构成基础解系；若
\(a_4=0\)，选择任意非零的 \(a_j\) 作主元即可。\(b=-s\) 时，基础解系为 \[
\left\{\begin{pmatrix}1\\1\\1\\1\end{pmatrix}\right\}.
\]

\subsection{四、幂等矩阵}\label{ux56dbux5e42ux7b49ux77e9ux9635}

（1）由 \(A(A-E_n)=0\)，有 \(r(A)+r(A-E_n)\le n\)；另一方面 \[
n=r(E_n)=r(A-(A-E_n))\le r(A)+r(A-E_n).
\] 故两者之和等于 \(n\)。

（2）\(A^2-A=0\)，故任一特征值 \(\lambda\) 满足
\(\lambda(\lambda-1)=0\)，即 \(\lambda=0\) 或
\(1\)。对应特征子空间的维数之和为 \[
n-r(A)+n-r(A-E_n)=2n-n=n.
\] 因此 \(A\) 可对角化，且相似于
\(\operatorname{diag}(E_r,0)\)。相似保持秩，而该对角矩阵的秩是
\(r\)，所以 \(r=r(A)\)。

（3）由（2），存在可逆 \(U\) 使 \[
U^{-1}AU=\begin{pmatrix}E_r&0\\0&0\end{pmatrix}.
\] 由 \(AB=BA\) 得 \[
(U^{-1}AU)(U^{-1}BU)=(U^{-1}BU)(U^{-1}AU).
\] 比较分块矩阵，得到 \[
U^{-1}BU=\begin{pmatrix}B_1&0\\0&B_2\end{pmatrix},
\] 其中 \(B_1,B_2\) 分别为 \(r\) 阶、\(n-r\) 阶矩阵。\(B^2=B\) 保证
\(B_1^2=B_1,B_2^2=B_2\)。再由（2），存在相应阶数的可逆矩阵 \(V_1,V_2\)
将二者对角化。取 \[
V=\begin{pmatrix}V_1&0\\0&V_2\end{pmatrix},
\] 则 \[
\begin{aligned}
V^{-1}U^{-1}BUV&=\begin{pmatrix}V_1^{-1}B_1V_1&0\\0&V_2^{-1}B_2V_2\end{pmatrix},\\
V^{-1}U^{-1}AUV&=\begin{pmatrix}V_1^{-1}E_rV_1&0\\0&V_2^{-1}0V_2\end{pmatrix}
=\begin{pmatrix}E_r&0\\0&0\end{pmatrix}.
\end{aligned}
\] 两式均为对角矩阵。取 \(P=UV\) 即得结论。

\subsection{五、相似与合同}\label{ux4e94ux76f8ux4f3cux4e0eux5408ux540c}

\subsubsection{1.
相似保持特征值，逆命题不成立}\label{ux76f8ux4f3cux4fddux6301ux7279ux5f81ux503cux9006ux547dux9898ux4e0dux6210ux7acb}

存在可逆矩阵 \(P\) 使 \(P^{-1}AP=B\)，于是 \[
\begin{aligned}
|xE-B|&=|xE-P^{-1}AP|\\
&=|P^{-1}(xE-A)P|\\
&=|P|^{-1}|xE-A|\,|P|=|xE-A|.
\end{aligned}
\] 所以相似矩阵有相同的特征多项式，因而有相同特征值。

反例为 \[
A=\begin{pmatrix}0&0\\0&0\end{pmatrix},\qquad
B=\begin{pmatrix}0&1\\0&0\end{pmatrix}.
\] 二者特征值均为零，但秩不同，故不相似。

\subsubsection{2.
给定矩阵不相似，但合同}\label{ux7ed9ux5b9aux77e9ux9635ux4e0dux76f8ux4f3cux4f46ux5408ux540c}

\[
A=3E_3-\begin{pmatrix}1\\1\\1\end{pmatrix}(1,1,1),
\] 由秩一矩阵性质，特征值为 \(3,3,3-(1,1,1)(1,1,1)^T=0\)。\(B\)
的非零对角块为 \[
\begin{pmatrix}2&1\\1&2\end{pmatrix}
=E_2+\begin{pmatrix}1\\1\end{pmatrix}(1,1),
\] 其特征值为 \(1,1+(1,1)(1,1)^T=3\)，故 \(B\) 的特征值为
\(3,1,0\)。两个矩阵的特征值不同，所以不相似；二者均为实对称矩阵，正、负惯性指数分别均为
\(2,0\)，故合同。

\subsection{六、线性变换的矩阵、核与值域}\label{ux516dux7ebfux6027ux53d8ux6362ux7684ux77e9ux9635ux6838ux4e0eux503cux57df}

（1）由题给关系有 \[
\varepsilon_4=\eta_4,\quad\varepsilon_3=\eta_3,\quad
\varepsilon_2=\eta_2+\eta_3,\quad
\varepsilon_1=\eta_1-\eta_2-\eta_3.
\] 以 \(\eta_j\) 的旧基坐标作为列，过渡矩阵及其逆为 \[
P=\begin{pmatrix}1&0&0&0\\1&1&0&0\\0&-1&1&0\\0&0&0&1\end{pmatrix},\qquad
P^{-1}=\begin{pmatrix}1&0&0&0\\-1&1&0&0\\-1&1&1&0\\0&0&0&1\end{pmatrix}.
\] 因此新基矩阵为 \[
\begin{aligned}
P^{-1}AP
&=\begin{pmatrix}1&2&3&-6\\2&0&-2&0\\4&3&-1&6\\1&1&1&-3\end{pmatrix}
\begin{pmatrix}1&0&0&0\\1&1&0&0\\0&-1&1&0\\0&0&0&1\end{pmatrix}\\
&=\begin{pmatrix}3&-1&3&-6\\2&2&-2&0\\7&4&-1&6\\2&0&1&-3\end{pmatrix}.
\end{aligned}
\] 原OCR中乘积的中间矩阵末行缺少一个 \(1\)，以上由矩阵乘法复核恢复。

（2）行变换过程为 \[
\begin{aligned}
A&\longrightarrow\begin{pmatrix}1&2&3&-6\\0&-4&-8&12\\0&-1&-5&18\\0&-1&-2&3\end{pmatrix}\\
&\longrightarrow\begin{pmatrix}1&0&-1&0\\0&1&2&-3\\0&0&-3&15\\0&0&0&0\end{pmatrix}\\
&\longrightarrow\begin{pmatrix}1&0&0&-5\\0&1&0&7\\0&0&1&-5\\0&0&0&0\end{pmatrix}.
\end{aligned}
\] 所以 \(A\) 的前三列是极大线性无关组，\(\operatorname{Im}\sigma\)
的一组基为 \[
\begin{aligned}
\sigma(\varepsilon_1)&=\varepsilon_1+3\varepsilon_2+2\varepsilon_3+\varepsilon_4,\\
\sigma(\varepsilon_2)&=2\varepsilon_1+2\varepsilon_2+3\varepsilon_3+\varepsilon_4,\\
\sigma(\varepsilon_3)&=3\varepsilon_1+\varepsilon_2+\varepsilon_3+\varepsilon_4.
\end{aligned}
\] \(Ax=0\) 的通解为 \[
x=x_4\begin{pmatrix}5\\-7\\5\\1\end{pmatrix},
\] 故 \(\ker\sigma\) 的一组基为
\(5\varepsilon_1-7\varepsilon_2+5\varepsilon_3+\varepsilon_4\)。

\subsection{七、二次型}\label{ux4e03ux4e8cux6b21ux578b}

（1）由特征值之和得 \(a+b=\operatorname{tr}A=1-2+4=3\)，所以 \[
b=3-a,\qquad A=\begin{pmatrix}3-a&-2&0\\-2&a&-2\\0&-2&0\end{pmatrix}.
\] 由于 \(1,-2\) 都是特征值， \[
\begin{aligned}
0=|A-E|&=\begin{vmatrix}2-a&-2&0\\-2&a-1&-2\\0&-2&-1\end{vmatrix}\\
&=(a-1)(a-2)+4-4(2-a)=(a-1)(a+2),\\
0=|A+2E|&=\begin{vmatrix}5-a&-2&0\\-2&a+2&-2\\0&-2&2\end{vmatrix}\\
&=2(5-a)(a+2)-8-4(5-a)=-2(a-1)(a-4).
\end{aligned}
\] 联立得 \(a=1,b=2\)。

（2）此时 \[
A=\begin{pmatrix}2&-2&0\\-2&1&-2\\0&-2&0\end{pmatrix}.
\] 对特征值 \(4\)， \[
\begin{aligned}
A-4E&=\begin{pmatrix}-2&-2&0\\-2&-3&-2\\0&-2&-4\end{pmatrix}\\
&\longrightarrow\begin{pmatrix}1&1&0\\0&-1&-2\\0&-2&-4\end{pmatrix}
\longrightarrow\begin{pmatrix}1&0&-2\\0&1&2\\0&0&0\end{pmatrix}.
\end{aligned}
\] 取特征向量
\(v_1=(2,-2,1)^T\)。原OCR末尾额外的三阶矩阵不能与前面的行变换相容，此处以原矩阵的实际行变换恢复，不将其当作另一条有效等式。

对特征值 \(1\)， \[
\begin{aligned}
A-E&=\begin{pmatrix}1&-2&0\\-2&0&-2\\0&-2&-1\end{pmatrix}\\
&\longrightarrow\begin{pmatrix}1&-2&0\\0&-4&-2\\0&-2&-1\end{pmatrix}
\longrightarrow\begin{pmatrix}1&0&1\\0&2&1\\0&0&0\end{pmatrix}.
\end{aligned}
\] 取特征向量 \(v_2=(2,1,-2)^T\)。

对特征值 \(-2\)， \[
\begin{aligned}
A+2E&=\begin{pmatrix}4&-2&0\\-2&3&-2\\0&-2&2\end{pmatrix}\\
&\longrightarrow\begin{pmatrix}2&-1&0\\0&2&-2\\0&-2&2\end{pmatrix}
\longrightarrow\begin{pmatrix}2&0&-1\\0&1&-1\\0&0&0\end{pmatrix}.
\end{aligned}
\] 取特征向量 \(v_3=(1,2,2)^T\)。这三个向量两两正交，长度均为 \(3\)，故
\[
Q=\frac13\begin{pmatrix}2&2&1\\-2&1&2\\1&-2&2\end{pmatrix},\qquad Q^TQ=E.
\] 用 \(X=QY\)，得到标准形 \(f=4y_1^2+y_2^2-2y_3^2\)。原OCR的 \(Q\)
仅剩一行，以上根据三个已给特征向量恢复完整矩阵。

（3）\(A+tE\) 仍为实对称矩阵，特征值为 \(t+1,t-2,t+4\)。故当且仅当
\(t>2\) 时正定，当且仅当 \(t\ge2\) 时半正定。
\end{document}
